%% kapa: exam=1 round=27 year=2016 period=1 form=A subject=real_estate subject_ko=부동산학원론 seq=9 number=89
%% source: Q-Net 감정평가사 기출문제 (cst00309) artlSeq=5169402 file=extracted/1008/5169402/제27회 감정평가사 1차 문제/제27회 감정평가사 1교시문제(A형).pdf
%% conversion: kordoc 4.24.51 -> tools/gen1.py (auto); review=proofread-sonnet; flags=none
\begin{question}{89}
부동산경기변동에 관한 설명으로 옳지 않은 것은?
\begin{choices}
  \item 계절적 변동은 예기치 못한 사태로 초래되는 비순환적 경기변동 현상을 말한다.
  \item 부동산경기변동이란 일반적으로 상승과 하강국면이 반복되는 현상을 말한다.
  \item 건축착공량과 부동산거래량은 부동산경기를 측정할 수 있는 지표로 활용될 수 있다.
  \item 하향시장 국면이 장기화되면 부동산 공실률 증가에 의한 임대료 감소 등의 이유로 부동산 소유자에게 부담이 될 수 있다.
  \item 회복시장은 일반적으로 경기가 하향을 멈추고 상승을 시작하는 국면이다.
\end{choices}
\end{question}
\answer{1}
